\setcounter{section}{52}
\section{Equipos departamentales}

Nombre completo: Equipos departamentales. Servidores. Medidas de seguridad para equipos departamentales y servidores. Centros de proceso de datos: diseño, implantación y gestión.

\localtableofcontents

\subsection{Equipos departamentales: servidores}

Según su propósito se pueden ordenar por tamaño

\begin{description}
    \item[Miniordenador] Ordenador de nivel medio, cálculos complejos y multiusuario, en desuso por el downsizing.
    \item[microordenador] Ordenador de sobremesa o portátil
    \item[workstation] Ordenador de gran rendimiento, dedicado a tareas de cómputo pesado y con hardware de gran capacidad.
    \item[superordenador] Ordenador de gran capacidad y potencia para cálculos complejpos y tareas especiales
    \item[Ordenador central / Mainframe] Dedicado al procesamiento de grandes cantidades de datos, prioriza entrada y salida frente a cómputo
    \item[Servidores] Término genérico para referirse a ordenadores que proveen servicios a través de una red.
\end{description}

Un servidor es un equipo informático especializado, encargado de proveer de servicios a redes de datos, o a otros ordenadores. Los principales aspectos a valorar son:

\begin{description}
    \item[Fiabilidad] Elementos redundantes que permitan el funcionamiento continuado frente a fallos de hardware
    \item[Facilidad de operación y mantenimiento] Como herramientas de gestión remota y monitorización
    \item[Eficiencia energética] Gran parte del coste de un centro de datos es la electricidad y la refrigeración
    \item[Facilidad de actualización o ampliación]
\end{description}

Algunos formatos de servidor son:
\begin{description}
    \item[Enracables] preparados para ser integrados en un \textbf{rack} o armario. Permiten modularidad.
    \item[Blade] Ordenador enrackable de una unidad.
    \item[Torre] Ordenador no preparado para integrarse en otros sistemas.
    \item[Sistemas convergentes] Agrupan en un único equipo físico potencia de cálculo, almacenamiento y comuinicaciones. PErmiten gestión y soporte desde un punto común.
    \item[Sistemas hiperconvergentes] Utilizan tecnologías de virtualización para ofrecer infraestructura definida por software \textbf{SDDC}
\end{description}

\subsection{Medidas de seguridad para equipos departamentales y servidores}

Se pueden agrupar por niveles (parecido al sistema OSI)

\subsubsection{Nivel Físico}
Se tiene en cuenta el control de acceso al equipo, redundancia de suministro eléctrico, de discos, de comunicaciones y de refrigeración. Hay herramoientas de monitorización a nivel físico como Nagios o Cacti.

\subsubsection{Nivel lógico}
Se segmenta la red en zonas de seguridad, evitando comprometer toda la red. Se separan zonas DMZ, VLAN para diferentes grupos.

Uso de gestores de ancho de banda, cortafuegos y firewall. 

Sistemas gestores de eventos de seguridad. Protección anti spam y antivirus.

Realización de auditorías.

Bastionado, políticas de actualización de sistemas, Backups y copias remotas.

Procedimientos reglados para administración, como altas y bajas, conexiones remotas, gestión de usuarios, políticas de contraseñas, y de mínimo privilegio.

Se recomienda comprobar el Esquema Nacional de Seguridad RD 311/2022 y las guías de seguridad CCN-STIC

\subsection{Centros de proceso de datos: diseño, implementación y gestión}

Pueden clasificarse según los servicios proporcionados.

\begin{description}
    \item[Corporativos] Proveen servicios a la propia organización.
    \item[Servicios gestionados] Se administran por un tercero en nombre de la empresa. Esta alquila los equipos y su infraestructura.
    \item[Colocación] Gestionados por una empresa que alquila el espacio e infraestructura, para que otras empresas alojen su hardware.
    \item[Nube] Equipos y sistemas virtualizados como servicio.
    \item[Principal] Provee de servicios de forma habitual. Proveyendo de forma continuada y parando solo por intervención o desastre.
    \item[De respaldo] destinado a la recuperación ante desastres, contingencias o como mecanismo alternativo en caso de problemas con el principal.
    \begin{description}
        \item[Sala fría] centro externo con toda la infraestructura para replicar el centro principal a partir de copias de seguridad. Requiere tiempo para transportar los datos.
        \item[Sala caliente] Funcionan de forma análoga al principal. Replicando los datos de este y pudiendo funcionar solo con las últimas actualizaciones de los datos.
        \item[Centro espejo] Réplica en tiempo real del principal. Requiere de redes de alta velocidad. Es el modelo mas caro.
        \item[Backup mútuo] Dos entidades usan sus centros de datos para hacer backup recíproco.
    \end{description}
\end{description}

\subsubsection{Subsistemas}
Según el estandar TIA-942, la infraestructura de los datacenter se compone por cuatro subsistemas:

\begin{description}
    \item[Telecomunicaciones] Cableado de armarios y horizontal, acceso redundante, cuarto de entrada, área de distribución, backbone, elementos activos, alimentación redundante, patch panels, documentación.
    \item[Arquitectura] Selección de ubicación, construcción, protección ignífuga, barreras de vapor, techos y pisos, áreas de oficina, UPS, batería y generador, control de acceso, videovigilancia, operación de red.
    \item[Sistema Eléctrico] Número de accesos, puntos de fallo, cargas críticas, redundancia y topología de UPS. puesta a tierra, apagado de emergencia, sistemas de corte, monitorización generadores.
    \item[Sistema mecánico] Climatización, presión positiva, canalización, refrigeración, condensadores, detección de incendios y extinción.
\end{description}

\subsubsection{Elementos}

\begin{description}
    \item[Centro de operaciones] Oficina donde se encuentran operadores y técnicos. Dispone de sistemas de monitorización e inspección remota. También se controla el acceso a las salas
    \item[Sala de entrada] Contiene los elementos de comunicaciones de proveedores de acceso, así como almacenes, carga y descarga, etc
    \item[Sala pricipal] Sala que contiene los servidores y otros equipos del CPD. Suele contar con un falso suelo para mantenimiento sencillo del cableado, los equipos se distribuyen en armarios en hileras.
    \item[Armarios de comiunicaciones] Se sigue una topología en estrella. Con un backbone que agrupa el grueso de las comunicaciones distribuyéndose a las distintas zonas.
\end{description}

\subsubsection{Diseño}
    Los elementos de diseño depende de la clasificación con la que se evalue. La disponibilidad se da en porcentaje de tiempo (al año), se adjunta equivalencia en horas:

    \begin{description}
        \item[Tier 1 Básico] no hay redundancia, se puede interrumpir el servicio programado o no, tiene una tasa  de 99.671\% ~ 28.82h
        \item[Tier 2 Redundante] Tiene componentes redundantes, suelos elevados y SAI, tasa de 99.741\% ~ 22.69h
        \item[Tier 3 Concurrentemente mantenible] Permite planificar mantenimiento sin afectar al servicio, pueden existir paradas no planificadas. Multiples líneas disponibles, pero una activa. Tasa de 99.982\% 1.58h
        \item[Tier 4 Tolerante a fallos] Sin pérdida de servicio por paradas programadas, capaz de soportar eventos no programados sin pérdida de servicio. Múltiples líneas de electricidad y refrigeración con múltiples componentes redundantes. Tasa de 99.995\% ~ 26m
    \end{description}

    En cuanto al sistema eléctrico, se recomienda considerar la red eléctrica como un suministro auxiliar de bajo coste, y la generación local como principal.

    Aquí depende de los grupos electrógenos, generadores a gasolina o diesel. Después los \textbf{Sistemas de Alimentación Ininterrumpida} y por último la distribución y cuadros eléctricos. Siendo recomendable 

    Para el sistema de refrigeración se intenta que las temperaturas estén entre 18C y 27C con una humedad de 30\%~50\%. Pudiendo usar distíntos sistemas de refrigeración como expansión directa, condensación en torre, pasillo frío/caliente, aire exterior\dots

    En cuanto a los sistemas de incendios, hay sistemas de detección: de calor, de humo, de gases, de llama\dots y de extinción: agua nebulizada o aspersión de gas inerte con distintos gases.

\subsubsection{Métricas de Eficiencia}
    \begin{description}
        \item[Power Usage Effectiveness (PUE)] Hace referencia a la eficacia de uso de la energía. Es la tasa de consumo total contra la del equipo TI. Un PUE de 1 es una asíntota. Se intenta bajar de 2, los datacenter nuevos bajan a 1.2 y algunos llegan a 1.06.
        \item[Carbon Usage Eggecrtiveness (CUE)] Hace referencia a las emisiones de carbono. Tiene un valor ideal de 0 (se compensa el carbono emitido)
        \item[Water Usage Effectiveness (WUE)] Efectividad del uso del agua en litros por Rendimiento obtenido. La mertrica del rendimiento es subjetiva a cada centro.
    \end{description}